\section{Sparse features and their evaluation}\label{sec:C25:sparse}
Sparse autoencoders approximate an activation vector using a small number of learned dictionary features. A common objective combines reconstruction error with a sparsity constraint or penalty. Let $x$ be a model activation, $f(x)$ sparse feature activations, and $D$ a decoder dictionary. The reconstruction term $\|x-Df(x)\|^2$ measures approximation in activation space. It does not directly measure whether feature labels are correct or whether all behaviorally important structure is represented.

Anthropic's 2023 monosemanticity report applies dictionary learning to a small transformer and analyzes feature interpretation, intervention, splitting, and limitations~\cite{bricken2023monosemanticity}. OpenAI's 2024 sparse-autoencoder work studies scalable training and evaluation of reconstruction, sparsity, and feature quality~\cite{gao2024sae}. These are primary research contributions with different scopes. Neither establishes that a sparse representation is a unique, complete ontology of the model's concepts.

Feature interpretation must examine more than maximally activating examples. A feature can look clean at its strongest activations while responding to unrelated contexts at lower levels. Sampling only impressive examples introduces selection bias. Evaluation should include activation ranges, random examples, counterexamples, and downstream interventions. Automatic feature descriptions generated by another model require independent checking; a fluent label can rationalize a mixed feature.

Reconstruction and causal faithfulness can diverge. A dictionary can reconstruct common activation patterns well while missing rare behavior that matters for security or authorization. Conversely, a feature intervention can change behavior by moving computation off the model's natural activation manifold. The resulting effect is evidence about that intervention, not necessarily evidence that the feature normally causes the behavior in the same way.

For agent inspection, the relevant question is whether a feature-level explanation improves a concrete decision or detects a failure under held-out conditions. Compare it with simpler probes, input-output monitors, and executable checks at matched access and cost. A feature-based monitor should be tested after model updates, context compaction, and tool changes because its calibration may depend on the original distribution. Internal access can provide useful evidence while remaining an additional system component with its own drift and false-alarm behavior.

\subsection{The 2026 monitorability boundary}\label{sec:C25:monitorability2026}
The preceding circuit and dictionary methods are methodological foundations, not descriptions of the entire October 2026 inspection frontier. Contemporary oversight also uses monitors of reasoning traces, actions and tool results. These observations are complementary to internal activations: a readable trace is neither a complete causal account nor a proof that an accepted action satisfies its specification.

The August implicit-influence study compares matched nudges delivered as incidental contextual details or as explicit instructions to use and conceal the nudge~\cite{duzan2026implicit}. Across its tested task formats and extended-thinking models, behavioral influence can persist while trace-based detection weakens substantially in the implicit condition. Consequently, a monitor validated mainly on explicit concealment instructions can face a different observation problem in ordinary biased decision making. Its GPT-5-based judge does not establish that the influence is undetectable in principle; for Claude 4.5 runs the observed channel is an API summary rather than raw thinking. These are material access limits, not evidence that all reasoning traces are uninformative.

OpenAI's April 2026 release supplies datasets and metric code, but explicitly describes its public evaluation scaffold as illustrative and mocked, omitting restricted environments~\cite{openai2026monitorrelease}. The inspected \path{gmean_minimal_criterion} routine computes an effect-normalized detection bound with eligibility conditions; \path{bootstrapped_gmean_metric} separates eligibility selection from estimation and resamples within instance--arm strata. Its reported quantity must not be relabeled ordinary classifier accuracy or a direct measurement of complete causal faithfulness. Report the eligible-instance fraction alongside the conditional metric; an undefined case is not evidence of zero monitor error. The source revision is \texttt{d14f181a0623}; it was inspected, not executed~\cite{openai2026monitorcode}.

The September Astra system card supplies an additional current boundary: it distinguishes CoT-only, action-only and full-context monitors and reports reduced CoT monitorability relative to earlier models on its evaluations~\cite{openai2026astramonitor}. These are provider-reported configured results. Full-context access contains evidence absent from a trace-only monitor, so a comparison must hold observation scope fixed before attributing a difference to the monitor model. The existence of a full-context deployment monitor does not establish the same visibility for an application consuming an ordinary API response or a summarized reasoning field.

A discriminating experiment crosses actor checkpoint, monitor checkpoint, observation scope, context transformation and adversarial awareness. Freeze complete trajectories for the monitor comparison; separately rerun actors when studying how awareness changes behavior. Compressing a trajectory after the fact measures evidence loss to the monitor, whereas letting the actor reason under compacted context changes both behavior and available evidence. Preserve that distinction when interpreting the interaction. Evaluate sensitivity to implicit influence, false alarms on benign reasoning, alert delay relative to effect commitment, and compute cost. A monitor that identifies an error after an irreversible effect provides diagnostic evidence but has not demonstrated preventive control. An internal signal should become a deployment gate only after its held-out calibration and intervention consequences have been tested at the actual action boundary.
